% Section 2: Phase 1 -- faithful h(K)=1 lattices

Phase~1 implements the trivial-class-number degenerate case of Lemma~2.2
of \cite{OpenAI2026Proof}: an imaginary quadratic field $K = \Q(\sqrt{-d})$
with class number $h(K) = 1$, where the ideal-class pigeonhole reduces to
the identity. For these fields the ring of integers $\OK$ embeds
standardly in $\R^2$ as a rank-$2$ lattice, and unit distances between
lattice points can be counted exactly.

\subsection{Coverage}

We implement the three smallest $h = 1$ cases (the Baker--Heegner--Stark
list \(1, 2, 3, 7, 11, 19, 43, 67, 163\) is encoded in the module-level
constant \texttt{CLASS\_NUMBER\_ONE\_DISCRIMINANTS}; the first three are
implemented):

\begin{tabular}{lll}
\toprule
$d$ & $\OK$ & Standard embedding into $\R^2$ \\
\midrule
$1$ & $\Z[i]$ (Gaussian integers)     & $a + bi \mapsto (a, b)$ \\
$2$ & $\Z[\sqrt{-2}]$                 & $a + b\sqrt{-2} \mapsto (a, b\sqrt{2})$ \\
$3$ & $\Z[\omega]$, $\omega = \frac{-1 + \sqrt{-3}}{2}$ & $a + b\omega \mapsto (a - b/2,\; b\sqrt{3}/2)$ \\
\bottomrule
\end{tabular}

\medskip
For each $d$ and bound $b \geq 0$, the module enumerates all
$(2b+1)^2$ lattice points with integer coefficients in $[-b, b]$
and counts unit-distance pairs exactly via a numpy $O(n^2)$
broadcast.

\subsection{Verifications}

\begin{verification}\label{ver:gauss_3x3}
For $\Z[i]$ with $b = 1$ (the $3 \times 3$ Gaussian grid, $n = 9$),
the artifact computes exactly $12$ unit-distance pairs: three rows of two
horizontal edges plus three columns of two vertical edges. Pinned in
\tname{tests/test_imaginary_quadratic_lattice.py::test_gaussian_3x3_grid_has_12_unit_distances}.
\end{verification}

\begin{verification}\label{ver:gauss_7x7}
For $\Z[i]$ with $b = 3$ (the $7 \times 7$ Gaussian grid, $n = 49$),
the artifact computes exactly $2 \cdot 7 \cdot 6 = 84$ unit-distance
pairs.
\textit{Source:
\tname{src/erdos_ant/imaginary_quadratic_lattice.py},
class \tname{ImaginaryQuadraticLattice.count_unit_distance_pairs}.}
Pinned in \tname{src/erdos_ant/verify.py} as the assertion
\tname{phase1_gaussian_7x7_grid_has_expected_84_pairs}.
\end{verification}

\begin{verification}\label{ver:eisenstein_dominance}
For each $b \in \{1, 2, 3, 4\}$, the Eisenstein-lattice (triangular)
unit-distance count strictly exceeds the Gaussian-lattice count at the
same $b$. The triangular lattice has $6$ nearest neighbours per
interior point versus the square lattice's $4$. Pinned in
\tname{test_eisenstein_unit_distances_exceed_gaussian_for_same_bound}.
\end{verification}

\subsection{Non-claims}

Phase~1 does not improve the asymptotic Erd\H{o}s exponent. With
$[K : \Q] = 2$, the construction is the classical square or triangular
grid; in particular $\nu(n) \leq n^{1 + O(1/\log \log n)}$ holds and the
$\delta > 0$ excess is not present. The asymptotic result strictly
requires $[K : \Q] \to \infty$, which is the Phase~3 setting.

Phase~1 is included in the artifact for two reasons: (i) it is a
self-contained, exact reference against which the Phase~3 construction
can be sanity-checked at low parameter values, and (ii) it documents
the standard complex embedding of $\OK$ in $\R^2$ unambiguously, which
is the substrate of every later step.
